\centerline{\bf \Huge Задачи со сдвигом}

\begin{flushright}
        \scriptsize{}
\end{flushright}

\problem{1}
Каждый год комиссия по климату заявляет, что среднегодовая температура на Земле выше, чем два года назад, но ниже, чем пять лет назад. Сколько лет подряд это заявление может быть правдой?

\problem{2}
Несколько мудрецов выстроились в колонну. На всех были либо чёрные, либо белые колпаки. Оказалось, что среди любых 10 подряд идущих мудрецов поровну мудрецов с белыми и чёрными колпаками, а среди любых 12 подряд идущих - не поровну. Какое наибольшее количество мудрецов могло быть?

\problem{3}
На границе круглого участка стоит забор, в котором ровно 2023 доски. У Тома Сойера есть краски трёх цветов. Он взялся покрасить весь забор, соблюдая следующие условия: любые две доски, между которыми ровно 2 или ровно 5 досок, должны быть окрашены в разные цвета. Справится ли Том?

\problem{4}
В ряд расположены несколько фишек двух цветов, причём оба цвета присутствуют. Известно, что фишки, между которыми 6 или 12 фишек, одинаковы. Какое наибольшее число фишек может быть?

\problem{5}
Какое наибольшее количество различных целых чисел можно выписать в ряд так, чтобы сумма каждых 11 подряд идущих чисел равнялась 100 или 101 ?
